%% raw_exams.tex -- Practice exam (Fall 2026) + official solutions, and Sai's graded midterm
%% (Gradescope export, Summer 2026). Raw transcription: every question verbatim.
\section{Practice exam (Fall 2026) and Sai's graded midterm (last quarter)}
\ixh{model evaluation}\ixh{model comparison}

\subsection{Practice exam with official solutions}
\ixh{overfitting}\ixh{KNN}
{\color{figgray}\itshape [note: the official key prints the correct choice(s) of every multiple-choice question and all
solution text in red bold; here the correct choices are in {\color{keyc}\bfseries green bold} and the solution text follows
``Official solution''. The question prompts are identical in the exam and in the key.]}

%% ------------------------------------------------------------------ Q1
\begin{qbox}{PE Q1}
\ixh{KNN}\ixh{overfitting}\ixh{tuning parameter}
1.~A \ix{KNN} classifier has excellent performance on the \ix{training data} but noticeably worse performance on new data.
Which change would be most likely to help?
\begin{enumerate}[label=\alph*.]
\item Decrease $k$
\item {\color{keyc}\bfseries Increase $k$}
\item Give variables with larger variances more weight
\item Include an interaction term
\end{enumerate}
\sol A KNN model that is \ix{overfitting} is likely because the value of $k$ is small. Increasing $k$ will likely improve
performance. Decreasing $k$ will only make the problem worse. Weighting variables with larger variances more heavily
will worsen performance\ixh{standardization}. Including an \ix{interaction term} is not necessary (KNN is
nonparametric\ixh{non-parametric model}).
\end{qbox}

%% ------------------------------------------------------------------ Q2
\begin{qbox}{PE Q2}
\ixh{flexibility}
2.~A student says, ``The most flexible model is always preferable because it can represent more complicated
relationships.'' Explain what is wrong with this reasoning.
\sol This is not true. There are many cases where the most flexible model\ixh{flexibility} is not preferable, either
because the more flexible model is less interpretable\ixh{interpretability} and the goal is interpretation\ixh{inference},
or because the more flexible model can fit too closely to patterns in the training data that do not exist in the
population, leading to \ix{overfitting} (low \ix{training error} and high \ix{testing error}).
\end{qbox}

%% ------------------------------------------------------------------ Q3
\begin{qbox}{PE Q3}
\ixh{accuracy}\ixh{overfitting}
3.~You fit Models A, B, and C:\par
\begin{tabular}{@{}lll@{}}\toprule
\textbf{Model} & \textbf{Training accuracy} & \textbf{Testing accuracy}\\\midrule
\textbf{A} & 71\% & 70\%\\
\textbf{B} & 86\% & 82\%\\
\textbf{C} & 100\% & 68\%\\\bottomrule
\end{tabular}\par
Which model would you prefer based only on this information? What appears to be happening with Model C?
\sol You should prefer Model B; it produces the highest testing \ix{accuracy}\ixh{testing data}\ixh{model comparison}.
Model C appears to be \ix{overfitting}. (Model B shows some overfitting also, but less so than Model C.)
\end{qbox}

%% ------------------------------------------------------------------ Q4
\begin{qbox}{PE Q4}
\ixh{Pearson correlation}\ixh{interaction term}\ixh{training error}
4.~A student makes these three claims:
\begin{enumerate}[label=\alph*.]
\item A predictor with zero \ix{Pearson correlation} with the outcome cannot improve prediction.
\item Two highly correlated predictors necessarily interact.
\item A model with lower training error is always preferable.
\end{enumerate}
For each claim, state whether it is true or false and explain.
\sol For (a), this is false. Predictors with zero Pearson correlation with the outcome do not have a linear
\ix{correlation} with the outcome but may very well be able to improve predictions in nonlinear\ixh{nonlinear relationship}
and/or empirically driven models. For (b), this is also false. Interactions\ixh{interaction term} have nothing to do with
whether predictors are correlated\ixh{collinearity}. An interaction means the relationship between X1 and Y differs
depending on the value of X2. For (c), this is also false. Training error\ixh{training error} is not a good estimator of
\ix{testing error}. As training error goes to zero, testing error almost always increases.\ixh{overfitting}
\end{qbox}

%% ------------------------------------------------------------------ Q5
\begin{qbox}{PE Q5}
\ixh{KNN}\ixh{classification}
5.~A new observation has the following seven nearest neighbors, already ordered from closest to farthest:\par
\begin{tabular}{@{}|l|c|c|c|c|c|c|c|@{}}\hline
Neighbor & 1 & 2 & 3 & 4 & 5 & 6 & 7\\\hline
Class & A & B & B & A & A & B & B\\\hline
\end{tabular}\par
Predict the new observation's class using $k = 1$, 3, 5, and 7.
\sol Using $k = 1$, we predict A; using $k = 3$, we predict B; using $k = 5$, we predict A; using $k = 7$,
we predict B.\ixh{k-nearest neighbors}
\end{qbox}

%% ------------------------------------------------------------------ Q6
\begin{qbox}{PE Q6}
\ixh{standardization}\ixh{KNN}
6.~A student standardizes\ixh{standardization} age but forgets to standardize annual income before fitting KNN. Explain why
this is a problem even though both variables are numeric.
\sol This is a problem because \ix{KNN} is a distance-based algorithm\ixh{distance metric}\ixh{Euclidean distance}.
Income most likely has a \underline{much} larger variance than age, and therefore will have a much larger impact on
distance calculations if it is not also standardized\ixh{scaling}.
\end{qbox}

%% ------------------------------------------------------------------ Q7
\begin{qbox}{PE Q7}
\ixh{KNN}\ixh{target leakage}\ixh{identifier variable}\ixh{recipes}
7.~A researcher wants to use KNN to predict whether or not a bank customer will default on a loan. They write the
following code:
\begin{rcode}
loan_recipe <-
  recipe(defaulted ~ ., data = loans) |>                    (*@\ixch{recipe()}\ixch{|>}@*)
  step_dummy(all_nominal_predictors(), one_hot = TRUE)      (*@\ixch{step\_dummy()}\ixch{all\_nominal\_predictors()}\ixch{one\_hot}@*)

loan_model <-
  nearest_neighbor(neighbors = 5) |>                        (*@\ixch{nearest\_neighbor()}\ixch{neighbors}@*)
  set_engine("kknn") |>                                     (*@\ixch{set\_engine()}\ixch{kknn}@*)
  set_mode("classification")                                (*@\ixch{set\_mode()}@*)

loan_fit <-
  loan_model |>
  fit(defaulted ~ ., data = loans)                          (*@\ixch{fit()}@*)
\end{rcode}
Their data contain the following variables:\\
\texttt{application\_id}\\ \texttt{age}\\ \texttt{income}\\ \texttt{employment\_type}\\ \texttt{credit\_score}\\
\texttt{default\_date}\\ \texttt{defaulted}

Identify at least four problems with this analysis. For each problem, explain why it matters and how you would
correct it.
\sol Here are all problems with this analysis, along with why they matter and how to correct:
\begin{itemize}[label=-]
\item ``.'' syntax will include all variables, including \texttt{application\_id}, which is an identifier\ixh{identifier variable},
  not an actual predictor, and \texttt{default\_date}, which would be explicit \ix{target leakage}. Remove those predictors.
\item There is no \ixc{step\_normalize()}, which is important for KNN because it is distance-based\ixh{normalization}.
  Add that to the recipe.
\item The recipe is not actually being used; \texttt{loan\_recipe} is defined, but the model is fit directly to the raw
  data.\ixh{workflow()}\ixh{add\_recipe()}
\item Employment\_type therefore hasn't actually been dummy encoded\ixh{dummy variable} for the fitted model.
\end{itemize}
Note that one hot encoding\ixh{one-hot encoding} is not a problem for KNN; KNN has no regression \ix{intercept} creating
the linear dependency problem.\ixh{rank deficiency}
\end{qbox}

%% ------------------------------------------------------------------ Q8
\begin{qbox}{PE Q8}
\ixh{scatterplot}\ixh{Pearson correlation}\ixh{nonlinear relationship}
8.~A \ix{scatterplot} of X and Y shows a very strong U-shaped relationship. Their Pearson correlation is 0.03. Which
statement is best?
\begin{enumerate}[label=\alph*.]
\item X and Y are essentially unrelated
\item The scatterplot is likely wrong
\item {\color{keyc}\bfseries There is little linear correlation, despite a strong nonlinear relationship}
\item X will not be a useful variable for predicting Y
\item Both C and D
\end{enumerate}
\sol The predictors have a strong U-shaped (nonlinear) relationship\ixh{nonlinear relationship}, but little linear
\ix{correlation}. This does not mean X will not be useful for predicting Y; nor does it mean they are unrelated. It also
does not mean the scatterplot is wrong.
\end{qbox}

%% ------------------------------------------------------------------ Q9
\begin{qbox}{PE Q9}
\ixh{interaction term}\ixh{exploratory data analysis (EDA)}
9.~Which \ix{EDA} finding is evidence of an \textbf{interaction}\ixh{interaction term}?
\begin{enumerate}[label=\alph*.]
\item X1 and X2 have a Pearson correlation of 0.91
\item {\color{keyc}\bfseries The relationship between X1 and Y is positive for one level of X2 and negative for the other}
\item X1 is strongly right-skewed and X2 is strongly left-skewed
\item X1 and Y have a Pearson correlation of 0.80
\item The relationship between X1 and Y is the same across both levels of X2
\end{enumerate}
\sol B is the definition of an interaction\ixh{interaction term}. \ix{Pearson correlation} does not imply the existence of
an interaction. Neither does skew\ixh{skewness}.
\end{qbox}

%% ------------------------------------------------------------------ Q10-12 stem
For Questions 10--12, suppose a university wants to predict whether a student will graduate within four years, using
information available \textbf{at the end of the student's first year}.\ixh{feature engineering}\ixh{predictor}\par
\begin{tabular}{@{}ll@{}}\toprule
Variable & Description\\\midrule
StudentID & Unique identifier\\
HS\_GPA & High school GPA\\
FirstYearGPA & GPA at end of first year\\
Major & Major at end of first year\\
CreditsAttempted & First-year credits attempted\\
CreditsCompleted & First-year credits completed\\
Dorm & First-year residence hall name\\
GraduationDate & Date of graduation, if graduated\\
FinalGPA & GPA upon leaving the university\\
Graduated4Years & Outcome\\\bottomrule
\end{tabular}

%% ------------------------------------------------------------------ Q10
\begin{qbox}{PE Q10}
\ixh{target leakage}\ixh{identifier variable}\ixh{outcome}
10.~Which variables clearly should \textbf{not} be included as predictors in the model? Select all that apply.
\begin{enumerate}[label=\alph*.]
\item {\color{keyc}\bfseries StudentID}
\item HS\_GPA
\item FirstYearGPA
\item {\color{keyc}\bfseries GraduationDate}
\item {\color{keyc}\bfseries FinalGPA}
\item Major
\item {\color{keyc}\bfseries Graduated4Years}
\item CreditsAttempted
\item CreditsCompleted
\end{enumerate}
\sol A is an ID variable\ixh{identifier variable}. D and E are both variables that would not be known at the end of the
first year\ixh{target leakage}. G is the \ix{outcome}.
\end{qbox}

%% ------------------------------------------------------------------ Q11
\begin{qbox}{PE Q11}
\ixh{dummy variable}\ixh{categorical predictor}\ixh{logistic regression}
11.~Which predictors would require dummy encoding before fitting a \ix{logistic regression}? Select all that apply.
\begin{enumerate}[label=\alph*.]
\item HS\_GPA
\item FirstYearGPA
\item {\color{keyc}\bfseries Major}
\item CreditsCompleted
\item {\color{keyc}\bfseries Dorm}
\end{enumerate}
\sol These are categorical variables\ixh{categorical predictor}, which would require dummy encoding\ixh{dummy variable}.
\end{qbox}

%% ------------------------------------------------------------------ Q12
\begin{qbox}{PE Q12}
\ixh{target leakage}\ixh{feature engineering}
12.~A student creates a new variable, CompletionRate. CompletionRate is calculated as the ratio of CreditsCompleted to
CreditsAttempted. Would including CompletionRate as a predictor constitute \ix{target leakage}? Explain.
\sol No. CompletionRate is just the proportion of credits that a student attempted in their first year that they were
actually able to complete, and it is not leakage (since the outcome is graduation at the end of \underline{four} years).
One could argue, however, that CompletionRate, CreditsCompleted, and CreditsAttempted should not all be included as
predictors simultaneously\ixh{collinearity} (though that is not the question being asked).
\end{qbox}

%% ------------------------------------------------------------------ Q13
\begin{qbox}{PE Q13}
\ixh{one-hot encoding}\ixh{dummy variable}\ixh{rank deficiency}\ixh{linear regression}
13.~A student wants to predict ``salary'' using ``experience'' and ``department.'' They fit the following model:
\begin{rcode}
salary_recipe <- recipe(salary ~ experience + department,
                        data = employees) %>%           (*@\ixch{recipe()}\ixch{\%>\%}@*)
  step_dummy(department, one_hot = TRUE)                (*@\ixch{step\_dummy()}\ixch{one\_hot}@*)

salary_model <- linear_reg() %>%                        (*@\ixch{linear\_reg()}@*)
  set_engine("lm")                                      (*@\ixch{set\_engine()}\ixch{lm}@*)

salary_workflow <- workflow() %>%                       (*@\ixch{workflow()}@*)
  add_recipe(salary_recipe) %>%                         (*@\ixch{add\_recipe()}@*)
  add_model(salary_model)                               (*@\ixch{add\_model()}@*)

salary_fit <- fit(salary_workflow, data = employees)    (*@\ixch{fit()}@*)
\end{rcode}
Suppose ``department'' has three levels, ``Sales,'' ``IT,'' and ``HR.''

After fitting the model, R reports that the model matrix is rank-deficient\ixh{rank deficiency} and that one or more
coefficients cannot be uniquely identified.\ixh{coefficient}
\begin{enumerate}[label=\alph*.]
\item What caused the problem?\\
  \textbf{\color{keyc}Official solution.}~One-hot encoding\ixh{one-hot encoding}; this creates three indicator variables,
  one per level of ``department,'' which results in a linear dependency that OLS regression\ixh{least squares} cannot
  handle.
\item Why does ``one\_hot = TRUE'' create this problem for an ordinary linear regression model with an intercept?\\
  \textbf{\color{keyc}Official solution.}~Because the model contains an \ix{intercept}, which is itself a column of ones.
  Therefore, the design matrix\ixh{design matrix} is not full rank.
\item How should the code be changed?\\
  \textbf{\color{keyc}Official solution.}~Set ``\ixc{one\_hot} = FALSE.''
\item After making this change, how many indicator variables will be created for ``department''?\\
  \textbf{\color{keyc}Official solution.}~Two.\ixh{dummy variable}\ixh{reference level}
\end{enumerate}
\end{qbox}

%% ------------------------------------------------------------------ Q14
\begin{qbox}{PE Q14}
\ixh{logistic regression}\ixh{posterior probability}
14.~A \ix{logistic regression} model predicts whether a patient has pneumonia based on symptoms and clinical
measurements. For one patient, the model produces a predicted probability of 0.78. Which interpretation is correct?
\begin{enumerate}[label=\alph*.]
\item The model is 78\% accurate for this patient.
\item The patient definitely has pneumonia.
\item {\color{keyc}\bfseries The fitted model assigns this patient a probability of 0.78 of belonging to the pneumonia class.}
\item 78\% of the patient's symptoms indicate pneumonia.
\end{enumerate}
\sol The model cannot be ``78\% accurate'' for an individual patient\ixh{accuracy}. Nor does this predicted probability
tell us anything about the patient's symptoms. We cannot be absolutely certain the patient either does or doesn't have
pneumonia. C is the correct definition.
\end{qbox}

%% ------------------------------------------------------------------ Q15
\begin{qbox}{PE Q15}
\ixh{threshold}\ixh{classification}
15.~A model predicts whether borrowers will default on a loan. Four borrowers receive the following predicted
probabilities:\par
\begin{tabular}{@{}ll@{}}\toprule
\textbf{Borrower} & \textbf{Predicted Probability}\\\midrule
A & .31\\ B & .48\\ C & .63\\ D & .81\\\bottomrule
\end{tabular}
\begin{enumerate}[label=\alph*.]
\item Using a \ix{threshold} of 0.50, which borrowers would be classified as likely to default?\\
  \textbf{\color{keyc}Official solution.}~Borrowers C and D.
\item Using a threshold of 0.70, which borrowers would be classified as likely to default?\\
  \textbf{\color{keyc}Official solution.}~Only Borrower D.
\item Explain what happens to the number of positive classifications when the threshold is increased.\\
  \textbf{\color{keyc}Official solution.}~When the threshold is increased, the number of positive classifications will
  tend to decrease.\ixh{true positive}\ixh{false positive}
\end{enumerate}
\end{qbox}

%% ------------------------------------------------------------------ Q16
\begin{qbox}{PE Q16}
\ixh{logistic regression}\ixh{coefficient}\ixh{log odds}
16.~A logistic regression predicts whether a customer will cancel a subscription. The \ix{coefficient} for the predictor
MonthlyFee is 0.04.

A student writes: ``A \$1 increase in the monthly fee causes a customer's probability of canceling to increase by 4\%.''

Identify \textbf{one problem} with this interpretation and provide a better interpretation.
\sol This statement implies that there is a linear relationship between monthly fee and probability of cancellation.
There is not. There is a linear relationship between monthly fee and \underline{log odds of cancellation}\ixh{log odds}, or
\ix{logit}. Also, we cannot assume that we can draw causal conclusions\ixh{association vs. causation}\ixh{confounding}. A
better interpretation would be either ``A \$1 increase in the monthly fee is \underline{associated} with a 0.04 increase in
\underline{log odds},'' or ``An increase in the monthly fee is associated with an increased probability of cancellation''
(this is known because the sign of the coefficient is positive).
\end{qbox}

%% ------------------------------------------------------------------ Q17
\begin{qbox}{PE Q17}
\ixh{confusion matrix}\ixh{accuracy}\ixh{TPR}\ixh{FPR}
17.~A model produces the following \ix{confusion matrix}:\par
\begin{tabular}{@{}lll@{}}\toprule
 & \textbf{Actually Positive} & \textbf{Actually Negative}\\\midrule
\textbf{Predicted Positive} & 72 & 18\\
\textbf{Predicted Negative} & 28 & 182\\\bottomrule
\end{tabular}\par
Calculate accuracy, true positive rate, and false positive rate.
\sol Accuracy\ixh{accuracy} is $(72 + 182)/(72 + 182 + 18 + 28) = 254/300 = 84.67\%$ (or 0.8467). The true positive
rate\ixh{TPR} is $72/(72 + 28) = 72/100 = 72\%$ (or 0.72). The false positive rate\ixh{FPR} is $18/(18+182) = 9\%$ (or
0.09).\ixh{true positive}\ixh{false positive}\ixh{false negative}\ixh{true negative}
\end{qbox}

%% ------------------------------------------------------------------ Q18
\begin{qbox}{PE Q18}
\ixh{AUC}\ixh{threshold}\ixh{ROC curve}
18.~A student says: ``I changed my classification threshold from 0.50 to 0.30, so the ROC AUC\ixh{ROC AUC} of my fitted
model will change.'' Is this correct? Explain.
\sol This is not correct. The ROC AUC measures the area under the \ix{ROC curve}, which is a curve that plots true positive
rates and false positive rates \underline{across thresholds}\ixh{threshold}. Therefore, using a different threshold value to
classify observations will have no impact whatsoever on the area under the ROC curve\ixh{AUC}.
\end{qbox}

%% ------------------------------------------------------------------ Q19
\begin{qbox}{PE Q19}
\ixh{ROC curve}\ixh{false negative}\ixh{TPR}\ixh{threshold}
19.~Suppose the \ix{ROC curve} below represents a model used to screen patients for a very dangerous disease. A
\ix{false negative} is considered substantially worse than a \ix{false positive}. Of points A, B, and C, which threshold
would you be most inclined to choose?\par
\fig{ggplot-style plot titled ``ROC Curve''; x-axis ``False Positive Rate'' 0.0--1.0 (ticks every 0.2), y-axis ``True
Positive Rate'' 0.0--1.0; light grey grid. A solid black stepped curve rises steeply from (0,0) (TPR $\approx$0.2 already at
FPR $\approx$0.01), bends near (0.15, 0.69), and flattens to (1,1); a dashed diagonal $y=x$ is the no-skill line. Three
black dots with red labels lie on the curve: \textbf{A} $\approx$ (FPR 0.05, TPR 0.45), \textbf{B} $\approx$ (0.22, 0.74),
\textbf{C} $\approx$ (0.47, 0.92). Moving A$\to$B$\to$C = lower threshold, higher TPR and higher FPR.}
\begin{enumerate}[label=\alph*.]
\item Point A
\item Point B
\item {\color{keyc}\bfseries Point C}
\item All points are equally good
\item Not enough information to decide
\end{enumerate}
\sol The problem implies that we are only interested in maximizing our true positive rate\ixh{TPR} (and therefore,
minimizing our false negative rate), regardless of our false positive rate\ixh{FPR}. Therefore, we would prefer the
\ix{threshold} corresponding to Point C (despite a false positive rate of almost 50\%).
\end{qbox}

%% ------------------------------------------------------------------ Q20
\begin{qbox}{PE Q20}
\ixh{LDA}\ixh{QDA}\ixh{covariance matrix}\ixh{decision boundary}
20.~You have 50,000 observations, two predictors, and clear evidence that the covariance structures\ixh{covariance matrix}
differ substantially across levels of your outcome variable. Which would you expect to have the greater potential to
capture the true \ix{decision boundary}?
\begin{enumerate}[label=\alph*.]
\item LDA
\item {\color{keyc}\bfseries QDA}
\item Both would do equally well
\item Neither would do well because they are both parametric models
\end{enumerate}
\sol \ix{LDA} assumes that the covariance structure does not differ and therefore would not do well. \ix{QDA} would
likely excel despite estimating more predictors\ixh{variance}; $p$ is not large, and there are a huge number of
observations (so parametric models\ixh{parametric model} are fine).
\end{qbox}

%% ==================================================================== SAI
\subsection{Sai's midterm (graded, 20/26)}
\ixh{model evaluation}
{\color{figgray}\itshape [note: source = Gradescope export \texttt{submission\_422008762.pdf}. PDF pp.~1--3 = Gradescope
rubric summary; pp.~4--8 = the scanned exam booklet (printed questions + Sai's pencil answers); pp.~9--10 = blank
scratch pages. Gradescope numbers the graded items ``Question 1''--``Question 12'', the booklet numbers them 1--9 (Q8 has
parts a--d), so both are given (``Gradescope Q12 / Q9''). In each rubric list, $\checkmark$ marks the rubric item the grader
applied; {\color{keyc}green} = $-0$ pts, {\color{warnc}red} = deductions, as on Gradescope.]}

\textbf{Gradescope header (PDF p.~1).} Midterm Exam --- Graded. Student: Sai Kadiyala. Total Points: \textbf{20 / 26 pts}.

\begin{rawsrc}{Booklet cover sheet (PDF p.~4)}
\textbf{PSTAT 131 Midterm Exam} --- Summer 2026.\quad NAME: \textit{Sai Kadiyala} [handwritten] (\textit{First and Last;
Please Print Clearly!})\quad NetID: [handwritten; omitted here] (\textit{This is what comes before the @ucsb.edu in your
UCSB Email Address})\quad Your TA (circle one): \textbf{Chunyin} [circled] \quad \textbf{Mallory}

\textbf{Instructions}
\begin{enumerate}
\item \textbf{Do not open this exam until you are told. Do not remove this cover sheet.}
\item \textbf{Identification:} Fill in your full name and UCSB NetID in the spaces above.
\item \textbf{Marking Answers:} You must indicate your answers to \textbf{all questions} (multiple choice and free response)
  on this exam booklet.
\item \textbf{Personal Aids:} This is an open-book and open-note exam. Any and all \textbf{physical materials} may be used
  during this exam. \textbf{Absolutely no} electronic or digital tools may be used.
\item \textbf{Prohibited Aids:} Use of generative AI tools (e.g., ChatGPT, Gemini), online resources, or communication with
  other students is strictly prohibited.
\item \textbf{Electronic Devices:} Cell phones, smart glasses, earbuds, and other unauthorized electronics must be turned
  off, put away, and kept out of sight for the duration of the exam.
\item \textbf{Timekeeping:} If a smart watch is used for timekeeping, it must be in airplane mode with all wireless
  communication (Wi-Fi/Bluetooth) disabled.
\item \textbf{Final Answers:} Round answers to 3--4 decimal places when appropriate. Answers must be legible; if multiple
  answers are given, the first will be graded, except for multiple choice questions; unless otherwise specified, multiple
  answers marked for Multiple Choice questions will receive a score of zero.
\item \textbf{Restroom Breaks:} If you must leave the room, give your exam paper and phone to the proctor to hold onto while
  you are out.
\item \textbf{Time Limit:} Exam ends at 4:35pm sharp; all work must stop immediately when called.
\item \textbf{Submission:} Turn in your exam to a proctor upon completion or when time is called. Show photo ID.
\end{enumerate}
\textbf{Academic Integrity Pledge.} \textit{In signing my name below, I certify that all work on this exam is my own. I
affirm that I have not used any unauthorized aids, technology (including generative AI tools), or electronic devices, nor
have I done anything to create an unfair advantage for myself or any other course member. I have not shared and will not
share information with any other student(s).} \quad Student Signature: [signed]
\end{rawsrc}

\textbf{Multiple Choice Questions} (booklet pp.~2--3 = PDF pp.~5--6)

%% ------------------------------------------------------------------ Sai Q1
\begin{qbox}{Sai Q1 = Gradescope Question 1 --- 2 / 2 pts}
\ixh{flexibility}\ixh{polynomial regression}\ixh{interaction term}
1.~(2 points) Which change generally makes a statistical learning method \textbf{more flexible}\ixh{flexibility}?
\begin{enumerate}[label=(\Alph*)]
\item Increasing $k$ in \ix{KNN}
\item Removing predictors from a linear model
\item Adding polynomial or interaction terms to a regression model\ixh{polynomial regression}\ixh{interaction term}
\item Using LDA\ixh{LDA} instead of QDA\ixh{QDA}
\item All of the above
\end{enumerate}
\par\textbf{\color{dblue}Sai's answer.}~(C) [bubble filled].
\grader
\begin{itemize}[leftmargin=1.2em]
\item[$\checkmark$] {\color{keyc}\bfseries $-$0 pts} Correct answer is C. All other answers would all result in a \textbf{less}
  flexible model.
\item[] {\color{warnc}\bfseries $-$2 pts} A, B, D, or E
\end{itemize}
\end{qbox}

%% ------------------------------------------------------------------ Sai Q2
\begin{qbox}{Sai Q2 = Gradescope Question 2 --- 2 / 2 pts}
\ixh{KNN}\ixh{decision boundary}\ixh{non-parametric model}
2.~(2 points) Why can \ix{KNN} represent nonlinear decision boundaries\ixh{decision boundary}\ixh{nonlinear relationship}
even without explicitly including polynomial or interaction terms?
\begin{enumerate}[label=(\Alph*)]
\item It assumes predictors are normally distributed
\item It does not assume equal variances across levels of the outcome
\item Classification is determined locally by neighboring observations
\item It estimates a different regression coefficient for every observation
\item It automatically creates squared terms
\item Both A and B
\end{enumerate}
\par\textbf{\color{dblue}Sai's answer.}~(C) [bubble filled].
\grader
\begin{itemize}[leftmargin=1.2em]
\item[$\checkmark$] {\color{keyc}\bfseries $-$0 pts} Correct answer is C. KNN does not make any assumptions about normality or
  equality of variances, it also doesn't create squared terms.
\item[] {\color{warnc}\bfseries $-$1 pt} Partial credit for an answer of B; students could be remembering that this assumption
  allows QDA\ixh{QDA} to represent nonlinear decision boundaries
\item[] {\color{warnc}\bfseries $-$2 pts} A, D, E, or F
\end{itemize}
\end{qbox}

%% ------------------------------------------------------------------ Sai Q3
\begin{qbox}{Sai Q3 = Gradescope Question 3 --- 2 / 2 pts}
\ixh{ROC curve}\ixh{threshold}\ixh{AUC}
\fig{the same ROC curve as practice-exam Q19, printed in black and white without grid lines: title ``ROC Curve'',
x-axis ``False Positive Rate'' 0.0--1.0, y-axis ``True Positive Rate'' 0.0--1.0, dashed diagonal reference line, solid
curve with three labelled dots: \textbf{A} $\approx$ (FPR 0.05, TPR 0.45), \textbf{B} $\approx$ (0.22, 0.74),
\textbf{C} $\approx$ (0.47, 0.92).}\par
3.~(2 points) Moving along the \ix{ROC curve} from Point A toward Point C corresponds generally to:
\begin{enumerate}[label=(\Alph*)]
\item Increasing the classification threshold
\item Decreasing the classification \ix{threshold}
\item Increasing the AUC
\item Decreasing the \ix{AUC}
\end{enumerate}
\par\textbf{\color{dblue}Sai's answer.}~(B) [bubble filled]. {\color{figgray}\itshape [note: a thin horizontal line runs
through option (C) across the whole page --- a scanner artifact, not a mark.]}
\grader
\begin{itemize}[leftmargin=1.2em]
\item[$\checkmark$] {\color{keyc}\bfseries $-$0 pts} Correct answer is B. Moving from Point A to Point C represents an increase
  in both true positive rate\ixh{TPR} and false positive rate\ixh{FPR}; this corresponds to decreasing the threshold,
  resulting in more predicted positives overall.
\item[] {\color{warnc}\bfseries $-$1 pt} Partial credit for an answer of A (wrong direction, but still recognizing that the
  threshold would be changing).
\item[] {\color{warnc}\bfseries $-$2 pts} C or D.
\end{itemize}
\end{qbox}

%% ------------------------------------------------------------------ Sai Q4
\begin{qbox}{Sai Q4 = Gradescope Question 4 --- 1 / 2 pts}
\ixh{identifier variable}\ixh{predictor}\ixh{categorical predictor}
4.~(2 points) A hospital is developing a model to predict whether a patient will be readmitted within 30 days after
discharge. Which of the following variables would be most reasonable to include as a \ix{predictor}?
\begin{enumerate}[label=(\Alph*)]
\item \texttt{PatientID}, a randomly assigned unique identifier for each patient.\ixh{identifier variable}
\item \texttt{HospitalID}, which identifies which of 15 hospitals treated the patient. Hospitals differ in their patient
  populations.\ixh{categorical predictor}
\item \texttt{AdmissionID}, a unique identifier assigned to each hospital admission.
\item \texttt{RecordID}, a sequential identifier (1, 2, 3, \dots) assigned to each row of the dataset.
\item None of the above.
\end{enumerate}
\par\textbf{\color{dblue}Sai's answer.}~(E) [bubble filled]. {\color{figgray}\itshape [note: the short slanted pencil
strokes through several words on this page are stray marks, not answers.]}
\grader
\begin{itemize}[leftmargin=1.2em]
\item[] {\color{keyc}\bfseries $-$0 pts} Correct answer is Hospital ID. The other variables are specifically defined as unique
  identifiers\ixh{identifier variable}, which should not be included as predictors, while patient populations differ across
  hospital.
\item[$\checkmark$] {\color{warnc}\bfseries $-$1 pt} Partial credit for answering all variables.
\item[] {\color{warnc}\bfseries $-$2 pts} A, C, or D
\end{itemize}
\end{qbox}

%% ------------------------------------------------------------------ Sai Q5
\begin{qbox}{Sai Q5 = Gradescope Question 5 --- 2 / 2 pts}
\ixh{QDA}\ixh{LDA}\ixh{covariance matrix}\ixh{variance}
5.~(2 points) Suppose the population covariance matrices\ixh{covariance matrix} really are identical across all classes.
What advantage does \ix{QDA} gain from estimating a separate covariance matrix for every class?
\begin{enumerate}[label=(\Alph*)]
\item A closer fit to the population decision boundary
\item No inherent advantage from the extra covariance flexibility; it may simply increase estimation variance\ixh{variance}
\item Perfect classification
\item QDA, unlike LDA, is nonparametric\ixh{non-parametric model}
\end{enumerate}
\par\textbf{\color{dblue}Sai's answer.}~(B) [bubble filled].
\grader
\begin{itemize}[leftmargin=1.2em]
\item[$\checkmark$] {\color{keyc}\bfseries $-$0 pts} Correct answer is B; if the primary assumption of \ix{LDA} is being met,
  fitting QDA will not provide any advantage, and will only increase \ix{variance}.
\item[] {\color{warnc}\bfseries $-$2 pts} A, C, or D
\end{itemize}
\end{qbox}

%% ------------------------------------------------------------------ Sai Q6
\begin{qbox}{Sai Q6 = Gradescope Question 6 --- 0 / 2 pts}
\ixh{interaction term}\ixh{main effect}
6.~(2 points) Consider the following fitted values:\par
\begin{tabular}{@{}lcc@{}}\toprule
 & \textbf{Treatment A} & \textbf{Treatment B}\\\midrule
\textbf{Low dosage} & 20 & 25\\
\textbf{High dosage} & 30 & 50\\\bottomrule
\end{tabular}\par
Which feature of these results most strongly suggests an interaction\ixh{interaction term}?
\begin{enumerate}[label=(\Alph*)]
\item Treatment B always has the larger fitted value
\item Increasing dosage changes the fitted value by 10 under Treatment A but by 25 under Treatment B
\item All four fitted values are different
\item Increasing dosage increases the fitted value across both treatments
\item Dosage has two levels
\item Both B and D
\end{enumerate}
\par\textbf{\color{dblue}Sai's answer.}~(F) [bubble filled].
\grader
\begin{itemize}[leftmargin=1.2em]
\item[] {\color{keyc}\bfseries $-$0 pts} Correct answer is B. The definition of an interaction is that the relationship between
  one variable and the outcome \textbf{differs} across levels of the other variable. A change in dosage has more of an
  effect on Treatment B than it does on Treatment A; this is what suggests an interaction.
\item[$\checkmark$] {\color{warnc}\bfseries $-$2 pts} A, C, D, E, or F
\end{itemize}
\end{qbox}

%% ------------------------------------------------------------------ Sai Q7
\begin{qbox}{Sai Q7 = Gradescope Question 7 --- 0 / 2 pts}
\ixh{logistic regression}\ixh{coefficient}\ixh{threshold}\ixh{log odds}
7.~(2 points) In a \ix{logistic regression} model, the coefficient for \texttt{age} is positive. Increasing one customer's
age from 40 to 41 does not change that customer's predicted class. Why is this possible?
\begin{enumerate}[label=(\Alph*)]
\item A positive coefficient does not affect predicted probabilities
\item The predicted probability can increase without crossing the classification threshold
\item Logistic regression coefficients affect only training observations
\item Age must interact with another predictor before it can affect classification
\end{enumerate}
\par\textbf{\color{dblue}Sai's answer.}~(A) [bubble filled].
\grader
\begin{itemize}[leftmargin=1.2em]
\item[] {\color{keyc}\bfseries $-$0 pts} Correct answer is B. An increase in age of 1 unit will have a positive effect on the
  \ix{log odds} (and therefore on the predicted probability), but that will only result in a changing classification if it
  results in crossing the \ix{threshold}.
\item[$\checkmark$] {\color{warnc}\bfseries $-$2 pts} A, C, or D
\end{itemize}
\end{qbox}

\textbf{Free Response Questions} (booklet pp.~4--5 = PDF pp.~7--8)

%% ------------------------------------------------------------------ Sai Q8
\begin{qbox}{Sai Q8a--d = Gradescope Questions 8--11 --- 8 / 8 pts}
\ixh{threshold}\ixh{sensitivity}\ixh{specificity}\ixh{ROC curve}
8.~A bank uses a model to flag credit-card transactions as potentially fraudulent. The classification \ix{threshold}
previously used was 0.40; the bank's lead data scientist decides to increase the threshold to 0.70.

Explain the effect, if any, that this change would generally have on each of the following:

\textbf{(a) (2 points) Sensitivity}\ixh{sensitivity} \hfill [Gradescope Question 8 / Q8a --- \textbf{2 / 2 pts}]
\par\textbf{\color{dblue}Sai's answer.}~``Sensitivity would be lowered'' \quad[margin note, to the right:] ``as [?] high
sensitivity $\to$ choose a point with high TPR\ixh{TPR}, even if FPR is higher $\to$ this means lower threshold'' [curved
arrow back to] ``so we are doing the opposite here''
\grader
\begin{itemize}[leftmargin=1.2em]
\item[$\checkmark$] {\color{keyc}\bfseries $-$0 pts} Sensitivity is the true positive rate (TPR). Increasing the threshold makes it
  harder to classify observations as positive. This will result in a lower sensitivity.
\item[] {\color{warnc}\bfseries $-$2 pts} Incorrectly saying that sensitivity would increase or stay the same.
\item[] {\color{warnc}\bfseries $-$1 pt} Partial credit
\end{itemize}

\textbf{(b) (2 points) Specificity}\ixh{specificity} \hfill [Gradescope Question 9 / Q8b --- \textbf{2 / 2 pts}]
\par\textbf{\color{dblue}Sai's answer.}~``Being stricter would decrease FPR\ixh{FPR} and the formula for specificity is
$1-\mathrm{FPR}$ so if FPR gets smaller specificity would go up.''
\grader
\begin{itemize}[leftmargin=1.2em]
\item[$\checkmark$] {\color{keyc}\bfseries $-$0 pts} Specificity\ixh{specificity} is the true negative rate\ixh{true negative} (or 1 minus the
  false positive rate). Increasing the threshold will result in a higher specificity.
\item[] {\color{warnc}\bfseries $-$2 pts} Incorrectly saying that specificity would decrease or stay the same.
\item[] {\color{warnc}\bfseries $-$1 pt} Partial credit
\end{itemize}

\textbf{(c) (2 points) The model's ROC curve}\ixh{ROC curve} \hfill [Gradescope Question 10 / Q8c --- \textbf{2 / 2 pts}]
\par\textbf{\color{dblue}Sai's answer.}~``The model's curve wouldn't necessarily change, only the dot of the threshold
would. Only determines which point on the ROC curve you opeate [sic] at.''
\grader
\begin{itemize}[leftmargin=1.2em]
\item[$\checkmark$] {\color{keyc}\bfseries $-$0 pts} A change in threshold will have no effect on the \ix{ROC curve} at all. The
  ROC curve exists at the model level and is made up of many different thresholds' combinations of true and false positive
  rates.
\item[] {\color{warnc}\bfseries $-$2 pts} Any other answer.
\item[] {\color{warnc}\bfseries $-$1 pt} Partial credit
\end{itemize}

\textbf{(d) (2 points) The overall number of transactions flagged for investigation} \hfill [Gradescope Question 11 / Q8d
--- \textbf{2 / 2 pts}]
\par\textbf{\color{dblue}Sai's answer.}~``when the threshold is increased, it means we require more ``certainty'' (higher
predicted probability)-so fewer observations exceed the high bar, so we get fewer positive predictions, more negative
predictions.''\ixh{false positive}
\grader
\begin{itemize}[leftmargin=1.2em]
\item[$\checkmark$] {\color{keyc}\bfseries $-$0 pts} The number of transactions flagged for investigation will generally
  decrease as the threshold is raised.
\item[] {\color{warnc}\bfseries $-$2 pts} Any other answer.
\end{itemize}
\end{qbox}

%% ------------------------------------------------------------------ Sai Q9
\begin{qbox}{Sai Q9 = Gradescope Question 12 --- 3 / 4 pts}
\ixh{collinearity}\ixh{normalization}\ixh{linear regression}\ixh{feature engineering}
9.~(4 points) Consider the following recipe. The model's goal is to predict the sale price of a home.
\texttt{square\_feet} represents the house size measured in feet; \texttt{square\_meters} represents the house size measured
in meters:
\begin{rcode}
house_recipe <-
  recipe(sale_price ~ ., data = houses) %>%           (*@\ixch{recipe()}\ixch{\%>\%}@*)
  step_mutate(                                        (*@\ixch{step\_mutate()}@*)
    square_meters = square_feet * 0.092903
  ) %>%
  step_normalize(all_numeric_predictors())            (*@\ixch{step\_normalize()}\ixch{all\_numeric\_predictors()}@*)
\end{rcode}
A student argues that because both \texttt{square\_feet} and \texttt{square\_meters} are normalized\ixh{normalization}, it is
appropriate to keep both in an ordinary \ix{linear regression}. Is this reasoning correct or incorrect? Explain.
{\color{figgray}\itshape [note: on the scan a thin horizontal scanner line runs through ``measured in meters''.]}

\par\textbf{\color{dblue}Sai's answer.}~``In simple linear regression\ixh{simple linear regression} we have this formula:
$Y=\beta_0+\beta_1X+\epsilon$ $\leftarrow$ error/noise'' [arrows from the formula: $Y\to$ ``response variable'',
$\beta_0\to$ ``intercept''\ixh{intercept}, $\beta_1\to$ ``slope''\ixh{slope}; a faint erased word under ``simple''
is [illegible]].\\
``The point I'm getting to is that, I personally believe it is not appropriate because the resulting value would be the
same (if only one of the two is used). If you normalize both in terms of price your result would be same as
standardizing/normalizing\ixh{standardization} means $\mu=0$ $\sigma=1$ in a [small sketch of a bell curve] normaly [sic]
distributed curve. So I believe it isn't.''
\grader
\begin{itemize}[leftmargin=1.2em]
\item[] {\color{keyc}\bfseries $-$0 pts} Students should recognize that this results in \ix{collinearity}, regardless of
  normalization. One predictor will still be a direct function of the other predictor.
\item[] {\color{warnc}\bfseries $-$2 pts} Logic is partially correct.
\item[] {\color{warnc}\bfseries $-$4 pts} Logic is completely incorrect.
\item[$\checkmark$] {\color{warnc}\bfseries $-$1 pt} Correct answer (that the reasoning is wrong) but doesn't specifically
  mention collinearity
\end{itemize}
\end{qbox}
{\color{figgray}\itshape [PDF pp.~9--10: blank pages (no writing).]}
